\section{Mathematical Formulation and Decision-Theoretic Foundation}
\label{sec:theory}

In this section, we establish the mathematical foundation governing operational wind power dispatch and market settlement. We formulate the operational problem under asymmetric balancing costs, derive the exact integral representation of excess operational risk, decompose it into representation error and calibration components, and establish formal boundary counterexamples delineating the limits of scalar calibration.

\subsection{Operational Dispatch Problem Formulation}
Let $X \in \mathcal{X} \subset \mathbb{R}^{L \times D}$ denote historical multivariate telemetry, and let $Y \in [0, 1]$ denote normalized turbine active power output over the target delivery horizon. In operational power system balancing, transmission system operators evaluate declared bids under asymmetric pinball dispatch penalties \cite{weron2014,zhang2014wind}:
\begin{equation}
\ell(p, Y; r) = c_+ \max(p - Y, 0) + c_- \max(Y - p, 0),
\label{eq:pinball_loss}
\end{equation}
where $p \in [0, 1]$ is the declared dispatch commitment, $c_+$ is the shortfall unit penalty (under-generation where declared power exceeds physical delivery, $p > Y$, requiring expensive tertiary reserve activation), $c_-$ is the surplus unit penalty (over-generation, $Y > p$, inducing downward redispatch or curtailment), and $r = c_+ / c_-$ represents the cost asymmetry ratio. Normalizing by $c_-$, the asymmetric operational loss is:
\begin{equation}
\mathcal{L}_{\text{disp}}(p, Y) = r \max(p - Y, 0) + \max(Y - p, 0).
\label{eq:norm_disp_loss}
\end{equation}
The expected operational risk under policy $p(X)$ is $\mathcal{R}(p) = \mathbb{E}_{X, Y} [\mathcal{L}_{\text{disp}}(p(X), Y)]$.

\subsection{Lemma 1: Exact Integral Excess Operational Risk}
Lemma~\ref{lem:excess_risk} formalizes a classical excess risk analytical identity in statistical decision theory and quantile regression, grounded in the foundational formulations of Koenker \cite{koenker2005} and electricity market bidding principles established by Pinson et al. \cite{pinson2007trading}. The unconstrained Bayes-optimal dispatch decision minimizing $\mathcal{R}(p)$ is the conditional quantile:
\begin{equation}
q_\tau(X) = F_{Y|X}^{-1}(\tau), \quad \text{where } \tau = \frac{1}{r + 1}.
\label{eq:bayes_quantile}
\end{equation}
For an operational cost asymmetry ratio $r=2$, the critical quantile is $\tau = 1/3 \approx 0.333$. Grounded in classical decision theory and proper scoring rules \cite{koenker2005,morales2014,pinson2007trading}, we express the relationship between point forecast deviation and downstream dispatch cost through the following integral identity:

\begin{lemma}[Decision-Theoretic Excess Operational Risk \cite{koenker2005,morales2014,pinson2007trading}]
\label{lem:excess_risk}
For any dispatch decision rule $p(X) \in [0, 1]$, the excess operational risk relative to the Bayes-optimal quantile $q_\tau(X)$ satisfies the exact integral identity:
\begin{equation}
\begin{aligned}
\mathcal{R}(p) - \mathcal{R}(q_\tau) &= (r + 1) \, \mathbb{E}_X \left[ \int_{q_\tau(X)}^{p(X)} \Big( F_{Y|X}(u) - \tau \Big) \, du \right] \\
&\ge 0,
\end{aligned}
\label{eq:exact_integral}
\end{equation}
which is non-negative everywhere. Under standard regularity conditions where $F_{Y|X}$ is strictly increasing on its support such that $q_\tau(X)$ is uniquely defined, equality holds if and only if $p(X) = q_\tau(X)$ almost surely.
\end{lemma}

\begin{proof}
For a fixed feature vector $X$, the expected conditional loss at commitment $p$ is:
\begin{equation}
\mathbb{E}[ \mathcal{L}_{\text{disp}}(p, Y) \mid X ] = r \int_0^p (p - y) dF(y) + \int_p^1 (y - p) dF(y),
\end{equation}
where $F(y) = F_{Y|X}(y)$. Taking the partial derivative with respect to $p$:
\begin{equation}
\begin{split}
\frac{\partial}{\partial p} \mathbb{E}[ \mathcal{L}_{\text{disp}}(p, Y) \mid X ] &= r F(p) - (1 - F(p)) \\
&= (r + 1) F(p) - 1.
\end{split}
\end{equation}
Recalling that $\tau = 1 / (r+1)$, this gradient simplifies to:
\begin{equation}
\frac{\partial}{\partial p} \mathbb{E}[ \mathcal{L}_{\text{disp}}(p, Y) \mid X ] = (r + 1) \Big( F(p) - \tau \Big).
\end{equation}
By the Fundamental Theorem of Calculus, integrating the gradient from $q_\tau(X)$ to $p(X)$ yields:
\begin{equation}
\begin{split}
&\mathbb{E}[ \mathcal{L}_{\text{disp}}(p, Y) \mid X ] - \mathbb{E}[ \mathcal{L}_{\text{disp}}(q_\tau, Y) \mid X ] \\
&\quad = (r+1) \int_{q_\tau(X)}^{p(X)} \Big( F(u) - \tau \Big) du.
\end{split}
\end{equation}
Because $F(u)$ is monotonically non-decreasing, for any $u > q_\tau(X)$ we have $F(u) \ge \tau$, so the integrand is non-negative; for any $u < q_\tau(X)$ we have $F(u) \le \tau$, and integrating backwards from $q_\tau(X)$ to $p(X)$ also yields a positive integral. Taking the outer expectation over $\mathcal{X}$ completes the proof.
\end{proof}

\subsection{Two-Stage Decomposition and Decoupled Decision Calibration}
When $p(X)$ is in the neighborhood of $q_\tau(X)$, expanding $F_{Y|X}(u)$ via a first-order Taylor series around $q_\tau(X)$ yields:
\begin{equation}
F_{Y|X}(u) - \tau \approx f_{Y|X}(q_\tau(X)) \cdot (u - q_\tau(X)),
\end{equation}
which produces the local quadratic excess risk approximation:
\begin{equation}
\mathcal{R}(p) - \mathcal{R}(q_\tau) \approx \frac{r+1}{2} \, \mathbb{E}_X \big[ f_{Y|X}(q_\tau(X)) \cdot ( p(X) - q_\tau(X) )^2 \big].
\label{eq:quadratic_approx}
\end{equation}

Now consider the standard two-stage paradigm where a forecaster predicts $f_\theta(X)$ aiming for conditional expectation $\mu(X) = \mathbb{E}[Y \mid X]$, and the dispatch commitment is formed via a scalar quantile offset $\delta$: $p(X) = f_\theta(X) + \delta$. Expressing the prediction error as:
\begin{equation}
p(X) - q_\tau(X) = \big( f_\theta(X) - \mu(X) \big) + \big( \delta - \Delta_\tau(X) \big),
\end{equation}
where $\Delta_\tau(X) = q_\tau(X) - \mu(X)$ denotes the conditional quantile-mean shift. Under homoscedastic curvature $w_\tau = \frac{r+1}{2} f_{Y|X}(q_\tau)$, substituting into Eq.~\eqref{eq:quadratic_approx} decomposes excess operational risk into three distinct terms:
\begin{equation}
\begin{aligned}
\mathcal{R}(p) - \mathcal{R}(q_\tau) &\approx w_\tau \, \mathbb{E}_X \big[ (f_\theta(X) - \mu(X))^2 \big] \\
&\quad + w_\tau \, \mathbb{E}_X \big[ (\delta - \Delta_\tau(X))^2 \big] \\
&\quad + 2 w_\tau \, \mathbb{E}_X \big[ (f_\theta(X) - \mu(X))(\delta - \Delta_\tau(X)) \big],
\end{aligned}
\label{eq:three_terms}
\end{equation}
where the components represent \textbf{Term I} (the risk of the calibrated estimator under the empirical representation), \textbf{Term II} (Quantile Mismatch / Calibration Dispersion), and \textbf{Term III} (Cross-Coupling Interaction Component), respectively.

We explicitly emphasize that this three-term decomposition is decision-theoretic and mathematically distinct from the classical statistical bias--variance decomposition: Term~I represents the risk of the calibrated estimator under the empirical representation, measuring the expected squared prediction error of the forecaster $f_\theta(X)$ relative to the true conditional expectation $\mu(X)$, and is never equated with prediction variance or parameter-estimation variance.

\subsection{Decision-Theoretic Deductions from the Three-Term Decomposition}
The three-term decomposition in Eq.~\eqref{eq:three_terms} provides a structured lens for interpreting how point predictions and quantile offsets interact:

\begin{remark}[Deduction 1: Pseudo-Advantage Masking in Uncalibrated Regimes]
\label{rem:masking}
In uncalibrated forecasting regimes ($\delta=0$), Term~II ($w_\tau \mathbb{E}[\Delta_\tau(X)^2]$) is an irreducible data-dependent constant representing the dispersion between conditional mean and target quantile. A physics-regularized model with inferior conditional mean accuracy (larger Term~I) can paradoxically achieve lower uncalibrated operational risk if its structural bias $(f_\theta(X) - \mu(X))$ aligns with the negative quantile gap $\Delta_\tau(X) = q_\tau(X) - \mu(X) < 0$ (for $r=2, \tau=1/3$). This alignment drives the interaction Term~III ($2 w_\tau \mathbb{E}[(f_\theta(X) - \mu(X))(-\Delta_\tau(X))] < 0$) negative, accidentally cancelling out a portion of representation error. Upon post-hoc quantile calibration, $\delta$ absorbs this systematic offset, neutralizing Term~III and revealing that the apparent physical advantage was an artifact of uncalibrated evaluation.

\textit{Diagnostic Caveat}: Term~III serves as an analytical diagnostic framework for explaining how structural bias can mask forecasting error under uncalibrated evaluation, rather than an empirically estimated scalar quantity. In operational telemetry where the true conditional distribution $F_{Y|X}$ is unknown, isolating the exact point-by-point magnitude of Term~III is intractable without strong parametric assumptions.
\end{remark}

\begin{remark}[Deduction 2: Regularization Distortion under Well-Sampled Regimes]
\label{rem:exposure}
Optimal quantile calibration ($\delta^* = \arg\min_\delta \hat{\mathcal{R}}(f_\theta + \delta)$) removes the unconditional component of Term~II. In well-sampled regimes where unconstrained models already achieve low representation error ($\text{Term~I}$), injecting an uncalibrated physical loss penalty can act as an uncompensated structural distortion that inflates Term~I, providing an analytical mechanism consistent with the empirical degradation observed under fixed priors in well-sampled settings.
\end{remark}

\begin{remark}[Deduction 3: Inductive Guidance and Prior Sensitivity under Sample Restriction]
\label{rem:inductive_guidance}
In operating regimes with restricted effective sample size, unconstrained empirical estimation is prone to over-fitting and elevated representation error under the empirical data sample (manifested as an inflated Term~I). Under sufficiently strong penalty weight ($\lambda=1.0$), genuine physical constraints provide inductive regularizing guidance that suppresses non-physical parameter excursions, securing robust downside risk reduction. Conversely, when this mapping is misspecified or permuted (as in our scrambled-prior negative control), the rigid penalty enforces counter-physical gradient updates, precipitating severe performance degradation and demonstrating acute vulnerability to prior misspecification under strong regularization.
\end{remark}

\subsection{Formal Boundary Counterexamples (Possibility Results)}
To ensure theoretical precision and delineate the limits of scalar calibration, we establish two formal possibility counterexamples demonstrating that decoupled calibration dominance is conditional rather than absolute.

Under \textit{Constant Shift Invariance} (Counterexample~\ref{ex:constant_shift}), if a physical constraint modifies a forecaster purely by a uniform additive shift $f_{\text{phys}}(X) = f_{\text{uncons}}(X) + c$, the subsequent optimal scalar calibration absorbs $c$ completely ($\delta_{\text{phys}}^* = \delta_{\text{uncons}}^* - c$), resulting in identical dispatch decisions and zero excess risk difference ($\Delta \mathcal{R}_{\text{disp}} \equiv 0$) despite an apparent increase in prediction MSE. Under \textit{Spatial Heteroscedasticity} (Counterexample~\ref{ex:heteroscedasticity}) with input-dependent variance $\sigma(X)$, a global scalar calibration fails to compensate for localized variance shifts; here, an estimator that introduces structural bias and exhibits higher MSE can achieve strictly lower operational risk than the true conditional mean following scalar quantile offset. These existence results prove that while scalar calibration successfully decouples mean prediction from global quantile alignment, it does not guarantee Pareto dominance over locally adaptive or asymmetric loss training under severe conditional heteroscedasticity. Complete mathematical proofs for both counterexamples are provided in \ref{sec:appendix_proofs}.

\subsection{Two-Settlement Wholesale Accounting and Real-Time BSC Imbalance Netting}
In wholesale electricity markets, wind power producers operate within a two-settlement framework coupling forward contracts and physical real-time imbalance settlement \cite{morales2014,pinson2007trading}. Total commercial revenue is expressed as:
\begin{equation}
R_{\text{total}} = P_{\text{F}} \cdot q + P_{\text{imb}} \cdot (y - q),
\label{eq:two_settlement}
\end{equation}
where $P_{\text{F}}$ denotes the forward contract price (e.g., bilateral PPA or day-ahead exchange clearing price), $q$ is the forward-contracted or nominated energy position in the stylized model, $y$ is actual physical generation volume, and $P_{\text{imb}}$ is the system imbalance cash-out price. Under the UK Balancing and Settlement Code (BSC) single-price regime (post-P305), energy imbalance cashflows aggregate over discrete half-hour Settlement Periods (SP):
\begin{equation}
E_{\text{SP}} = \sum_{t \in \text{SP}} (y_t - q_t) P_{\text{rated}} \Delta t \quad [\text{MWh}], \quad \Delta t = \frac{1}{6}\text{ h},
\label{eq:bsc_netting}
\end{equation}
where $P_{\text{rated}}$ denotes total wind farm nameplate capacity in MW. Physical shortfall occurs when $E_{\text{SP}} < 0$, incurring settlement charges evaluated at the System Buy Price $\lambda_t^{\text{SBP}}$. The heavy-tailed price distribution (mean \pounds 44.27/MWh, peak \pounds 1,509.80/MWh) acts as a high-magnitude risk multiplier. When extreme price spikes coincide with wind ramp errors, the cash-out penalty relative to average market conditions is amplified by up to $34.1\times$ (reflecting the ratio of the annual peak SBP of \pounds 1{,}509.80/\text{MWh} to the mean SBP of \pounds 44.27/\text{MWh}). However, as demonstrated in Section~\ref{sec:alignment}, evaluating forecasters on isolated unhedged balancing cashflows $P_{\text{imb}} \cdot (y - q)$ without forward revenue $P_{\text{F}} \cdot q$ creates severe metric distortion, establishing that non-negative asymmetric decision loss serves as the primary ground-truth metric.
